\section{Variable coefficients}\label{ch:varcoef}
\sectionbrief
{a rectangle $\Omega=I\times J\subset\R^2$ and a solution $u$ of $iu_t=-u_{xx}+V\,u$, i.e.\ $L_0u=Vu$ with $L_0=i\partial_t+\partial_x^2$, where the potential $V(x)$ or $V(t,x)$ is unknown. Solutions are \emph{strong} (class SF): $u\in C^{1,2}(\Omega)$ and the equation holds almost everywhere.}
{in physics the coefficients are rarely constant: a Schr\"odinger particle feels a potential. Phases~1--5 identified constant coefficients. Here the unknown is a function, and the first question is whether a single solution can reveal it at all.}
{extends Section~\ref{ch:ident} from $\C$ to function coefficients. It supplies Step~3c of the extended algorithm (Section~\ref{ch:extalgo}). The proofs are Dim.~A36--A40 and B78.}

\subsection{Too much freedom: a huge class identifies nothing}

Suppose we allow coefficients in any algebra $\mathcal K$ of functions closed under differentiation (all smooth functions, all analytic functions). Then even one solution that never vanishes carries no information beyond itself.

\begin{theorem}[{Rank-one annihilator \led{VAR-N1, VAR-T1(a)} \proved}]\label{thm:varn1}
Let $u\in C^\infty(\Omega)$ have no zeros, with $a=u_t/u$ and $b=u_x/u$ in $\mathcal K$. Then the annihilator of $u$ among differential operators with coefficients in $\mathcal K$ is generated by $\partial_t-a$ and $\partial_x-b$, and the module $\mathcal K[\partial]u$ is isomorphic to $\mathcal K$ itself. For an affine family $L_\theta=L_0+\sum\theta_kT_k$ the compatible parameters are $\theta_0+\ker\Phi_u$, so $\theta$ is identified iff $\Phi_u(\theta)=\sum\theta_kT_ku$ is injective.
\dimref{VAR-N1}
\end{theorem}

\begin{intuition}
If the coefficient class is closed under differentiation, a zero-free $u$ can write each of its own derivatives as ``a coefficient times $u$'', so every law is a law about $u$ itself and nothing more. To identify anything the class of coefficients must be restricted in advance.
\end{intuition}

\begin{caution}[title=Zeros matter]
With zeros the statement fails: $u=t-x^2$ near $0$ satisfies no law $\partial_t+\alpha\partial_x+\beta$ with smooth coefficients, because at the origin it would give $1=0$.
\end{caution}

\subsection{A potential $V(x)$ is read from any non-zero solution}

For a time-independent potential the situation is the opposite: nothing can hide.

\begin{theorem}[{Identification of $V(x)$ \led{VAR-T1(b)} \proved}]\label{thm:vart1b}
Let $u\not\equiv0$ be a strong solution with potential $V\in L^\infty_{\rm loc}(J)$, and $X_0=\{x:u(t,x)=0\ \forall t\}$. Then $X_0$ is closed and discrete; if $L_0u=W(x)u$ for another $W$, then $W=V$ almost everywhere; and
\[
V(x)=\frac{iu_t+u_{xx}}{u}\qquad\text{at every }(t,x)\text{ with }u(t,x)\ne0 .
\]
\dimref{VAR-T1b}
\end{theorem}

\begin{intuition}
If $u$ vanished along infinitely many vertical lines accumulating at $x_0$, cutting $u$ off to the right of $x_0$ would give another strong solution that vanishes on a half-strip. Schr\"odinger waves have infinite speed, and a Carleman estimate (a weighted energy inequality) shows that a solution vanishing on a strip vanishes everywhere. So $u$ cannot have such lines, and the formula $V=(iu_t+u_{xx})/u$ holds on a set of full measure.
\end{intuition}

\begin{example}{The harmonic oscillator}
For $V=x^2$ and the Hermite functions $h_0,\dots,h_3$ the set $X_0$ is empty or $\{0\}$, and $(iu_t+u_{xx})/u$ returns $x^2$ exactly. The same formula recovers the P\"oschl--Teller and Airy potentials.
\end{example}

\subsection{A potential $V(t,x)$: silent strips}

\begin{theorem}[{Time-dependent potentials \led{VAR-T1(c)} \proved}]\label{thm:vart1c}
For $V\in C(\Omega)$, the potentials compatible with $u$ are $V+\{W:W=0\text{ on }\Omega\setminus Z(u)\}$, where $Z(u)=\{u=0\}$; so $V$ is identified iff $\mathrm{int}\,Z(u)=\emptyset$. If $\mathrm{int}\,Z(u)\ne\emptyset$ there is an open interval $T'\subset I$ with $u=0$ on $T'\times J$ (a \emph{silent strip}). Silent strips occur: Tychonoff's series gives a solution vanishing for $t\le0$ and non-zero for $t>0$. Hence ``every $u\not\equiv0$ identifies $V(t,x)$'' is false.
\dimref{VAR-T1c}
\end{theorem}

\begin{falsebox}{Data that switch on from the boundary}{\falsest}
On a bounded domain without boundary conditions, a solution can be identically zero up to a time and non-zero afterwards (Tychonoff's series). While it is silent every potential supported in the silent time is compatible with it. The failure is in the time direction: unique continuation holds in $x$ but not in $t$.
\end{falsebox}

\subsection{Stability and the phase}

\begin{theorem}[{Stability of the reconstruction \led{VAR-P3} \proved}]\label{thm:varp3}
On a rectangle $K$ with $|u_1|\ge\delta>0$: $\|V_1-V_2\|_\infty\le\delta^{-1}(1+\|V_2\|_\infty)\|u_1-u_2\|_{C^{1,2}(K)}$; the map is not continuous from $C^1$ to $L^\infty$ (an explicit oscillation); and under an a priori bound $M$ on $\|e_{tt}\|_\infty,\|e_{xxx}\|_\infty$ one has a H\"older modulus with sharp exponent $1/3$. It is a modulus of continuity, not a certificate.
\dimref{VAR-P3}
\end{theorem}

\begin{theorem}[{Unobserved phase: $V$ modulo constants \led{VAR-P4} \proved}]\label{thm:varp4}
If only the class $\{e^{i\theta(t)}u\}$ is observed, the compatible potentials are exactly $\{V+c:c\in\R\}$. For $P(\partial_t,L_q)u=0$ with $P$ unknown, $(P,q)$ is defined only modulo $q\to q+c$, $P(\tau,\ell)\to P(\tau,\ell-c)$: a redundancy of the parametrisation.
\dimref{VAR-P4}
\end{theorem}

\begin{proposition}[{Dirac $1+1$ with a potential \led{VAR-P5} \proved}]\label{prop:varp5}
Finite speed of propagation leaves regions that the data never visit: for $m=V=0$ a travelling bump $\psi=f(x-t)v_+$ vanishes for $x\le0.5$ and $x\ge1.7$, so every potential supported there is compatible. The matrix potential is determined at $x_0$ iff $\mathrm{span}_t\{\psi(t,x_0)\}=\C^2$.
\dimref{VAR-P5}
\end{proposition}

\subsection{Other results of this phase}

\begin{itemize}
\item \textbf{Algebra of differential operators with polynomial coefficients} (Weyl algebra): for a $D$-finite $u$ the annihilator is a principal left ideal and $u$ identifies iff its derivatives are independent over $K$; Airy is irreducible and every solution identifies; the oscillator is reducible exactly for $E\in\pm(2\N+1)$ \led{VAR-T2a, VAR-T6, VAR-E1, VAR-E2}.
\item \textbf{Dynamical Lie algebras}: $\dim\mathfrak g_V<\infty$ iff $V\in\mathrm{span}\{1,x,x^2\}$; the modified P\"oschl--Teller potential has an explicit evolution through a Darboux chain \led{VAR-P8}.
\item \textbf{Non-autonomous splitting}: $E^{\rm mult}=\iint[A_2(s_1),A_1(s_2)]+O(t^3)$; commuting at equal times is not enough \led{VAR-P9}. Dyson--Phillips converges for bounded time-dependent perturbations \led{VAR-P11}; Floquet logarithms are bounded but not differential \led{VAR-P12}.
\end{itemize}
These are list cards of the catalogue (Part~IV); the closures made in Phase~10 are in Section~\ref{ch:phase10}.

\begin{takeaway}
A function-valued coefficient is identified by one solution exactly when the solution cannot hide: for $V(x)$ it cannot (infinite speed and unique continuation); for $V(t,x)$ it can fall silent; for a huge class of coefficients, nothing is identified. Stability is of order two and the phase costs one constant.
\end{takeaway}
